\section{Passover, Signs, and the Shape of the Chapter}

\subsection{A Passover sign without a Last-Supper narrative}

John dates the feeding by saying that Passover was near (6:4). That notice is
not decorative. The crowd crosses water; bread is supplied in a deserted place;
Moses and manna become the explicit terms of the later dispute; and the people
identify Jesus as the expected prophet before trying to make him king. The
chapter places new-exodus expectations around Jesus while steadily transferring
attention from a gift to its giver.

The Fourth Gospel does not narrate the words over bread and wine at the Last
Supper. It narrates the foot washing, farewell discourse, and Passion. That
literary fact explains why Christian readers have long asked whether John 6
performs some of the theological work done by the institution narratives in
the Synoptics and 1 Corinthians. It does not prove that every clause is a
liturgical transcript. The discourse occurs in Capernaum before the Passion,
not in the upper room. Its future ``bread that I will give'' points toward a
gift not yet accomplished within the narrative.

The Passover setting permits several mutually reinforcing horizons: manna,
exodus, divine provision, the approaching death of Jesus, and the Church's
Eucharistic reception. It does not license a simple equation in which the five
barley loaves are already consecrated Eucharistic species. John calls the
feeding a \emph{sign}; a sign discloses more than itself without ceasing to be
the concrete work narrated.

\subsection{A six-part movement}

The chapter's argument can be followed without pretending it was composed as a
modern outline:

\begin{enumerate}
\item \textbf{Sign and attempted kingship (6:1--15).} Jesus tests Philip,
receives Andrew's inadequate inventory, gives thanks, feeds the crowd, and
withdraws when the sign is converted into a political program.
\item \textbf{Sovereignty over the sea (6:16--21).} The one sought for bread
comes to the disciples over the water and identifies himself. The episode keeps
divine identity, not food alone, at the chapter's center.
\item \textbf{Misread sign and true work (6:22--34).} The crowd seeks more
food. Jesus contrasts perishing food with food abiding to eternal life and
names faith in the sent one as God's work.
\item \textbf{Bread, descent, and faith (6:35--47).} ``I am the bread of life''
is unfolded through coming, believing, the Father's drawing, the Son's descent,
and resurrection on the last day.
\item \textbf{Flesh, blood, eating, and abiding (6:48--59).} The discourse
returns to bread, now naming the gift as Christ's flesh for the world and
joining eating and drinking to abiding and resurrection.
\item \textbf{Crisis among disciples (6:60--71).} Murmuring reaches the
disciple circle. Jesus invokes ascension and the life-giving Spirit; many leave,
Peter confesses, and Judas is named under judgment.
\end{enumerate}

The repeated future resurrection formula---``I will raise him up on the last
day''---binds the faith language of verses 39--40 and 44 to the eating language
of verse 54. That repetition gives believing and eating a shared soteriological
horizon. It weighs against treating the two movements as unrelated sermons,
without by itself proving that every occurrence of bread has one sacramental
referent.

\subsection{Murmuring and narrative diagnosis}

John uses \emph{gongyz\=o}, ``murmur,'' for the objections at 6:41 and again
for the disciples' reaction at 6:61. The same Greek word-family recurs densely
in the Septuagint's manna narrative (Ex 16:2, 7--12), so John's wording can
evoke Israel's wilderness complaints. That echo does not erase the differences
among the questions. The crowd asks how Jesus can have come down from heaven
when his parents are known (6:42), then how he can give his flesh to eat
(6:52). The disciples call the saying hard (6:60). The text does not say that
every questioner understood an identical claim. A responsible reading follows
each objection in sequence and does not use ``murmuring'' as permission to
caricature all difficulty as bad faith.

